\documentclass[11pt,a4paper,onecolumn]{article}
\usepackage[margin=2.5cm]{geometry}
\usepackage{fontspec}
\usepackage{amsmath,amssymb}
\usepackage{booktabs,longtable,array}
\usepackage{enumitem}
\usepackage[hidelinks]{hyperref}
\setlist{itemsep=2pt,topsep=4pt}

% AUTO-GENERATED by scripts/build_gwaya_v3_paper_tables.py from results/gwaya_v3_oracle_eval.json
\newcommand{\corpusSha}{ffc051e3058783ca}
\newcommand{\repeats}{5}
\newcommand{\leanVersion}{Lean 4.33.1}
\newcommand{\rustVersion}{rustc 1.97.1}
\newcommand{\nInScope}{26}
\newcommand{\nLimit}{2}
\newcommand{\nTotal}{28}
\newcommand{\allDet}{yes}
\newcommand{\confusionTable}{%
\begin{tabular}{lrrrrr}\toprule
Subset & $n$ & true accept & true reject & false accept & false reject\\\midrule
In scope & 26 & 9 & 17 & 0 & 0\\
Known blind spots & 2 & 0 & 0 & 2 & 0\\
\bottomrule\end{tabular}}
\newcommand{\latencyTable}{%
\begin{tabular}{lrrr}\toprule Gate & samples & median (ms) & p95 (ms)\\\midrule
python & 60 & 0.16 & 633.08\\
rust & 30 & 96.83 & 145.36\\
lean & 50 & 663.38 & 911.59\\
\bottomrule\end{tabular}}
\newcommand{\caseTable}{%
\begin{longtable}{p{0.30\linewidth}lllp{0.12\linewidth}}\toprule
Case & Category & Label & Gate & Correct\\\midrule\endhead
\texttt{py\_ok\_sum} & valid & accept & accept & yes\\
\texttt{py\_ok\_walrus} & valid & accept & accept & yes\\
\texttt{py\_stub\_pass} & stub & reject & reject & yes\\
\texttt{py\_stub\_ellipsis} & stub & reject & reject & yes\\
\texttt{py\_stub\_mock} & stub & reject & reject & yes\\
\texttt{py\_empty} & empty & reject & reject & yes\\
\texttt{py\_comment\_only} & empty & reject & reject & yes\\
\texttt{py\_docstring\_only} & empty & reject & reject & yes\\
\texttt{py\_syntax\_error} & syntax & reject & reject & yes\\
\texttt{py\_semantic\_bug} & limitation & reject & accept & \textbf{no}\\
\texttt{py\_spec\_passed} & test\_spec & accept & accept & yes\\
\texttt{py\_spec\_failed} & test\_spec & reject & reject & yes\\
\texttt{rs\_ok\_add} & valid & accept & accept & yes\\
\texttt{rs\_ok\_generic} & valid & accept & accept & yes\\
\texttt{rs\_borrow\_move} & type & reject & reject & yes\\
\texttt{rs\_type\_mismatch} & type & reject & reject & yes\\
\texttt{rs\_syntax} & syntax & reject & reject & yes\\
\texttt{rs\_semantic\_bug} & limitation & reject & accept & \textbf{no}\\
\texttt{lean\_ok\_comm} & valid & accept & accept & yes\\
\texttt{lean\_ok\_simp} & valid & accept & accept & yes\\
\texttt{lean\_ok\_decide} & valid & accept & accept & yes\\
\texttt{lean\_false} & type & reject & reject & yes\\
\texttt{lean\_sorry} & escape & reject & reject & yes\\
\texttt{lean\_axiom} & escape & reject & reject & yes\\
\texttt{lean\_native\_decide} & escape & reject & reject & yes\\
\texttt{lean\_axiom\_lexical\_bypass} & escape & reject & reject & yes\\
\texttt{lean\_example\_bypass} & escape & reject & reject & yes\\
\texttt{lean\_comment\_fp} & regression & accept & accept & yes\\
\bottomrule\end{longtable}}


\title{GWAYA v3: A Fail-Closed Verification Gate for LLM-Generated Code\\[4pt]
\large A scope-limited report restricted to executable, reproducible evidence}
\author{Xavier Callens\\AutoevolveAI / ANSE open-source project}
\date{3 October 2026 \quad (version 3.1.0)}

\begin{document}
\maketitle

\begin{abstract}
GWAYA is a verification layer that sits between a code-generating language model
and its user. Version~3 makes the layer \emph{fail-closed}: a candidate is
accepted only if a deterministic check positively succeeds, and every situation
in which no check can run (missing toolchain, missing generator, offline
fallback) yields \emph{unverified} rather than approval. The gate combines an
AST-level stub audit, a Python parser with an isolated sandboxed test-execution
oracle, the \texttt{rustc} type and borrow checker, and the Lean~4 kernel
together with a \texttt{\#print axioms} audit and dynamic toolchain auto-discovery.
We report only what is backed by tests in the repository: on an author-constructed
labelled corpus of \nInScope{} in-scope candidates the gate made no errors, it was
deterministic over \repeats{} repetitions, and on \nLimit{} deliberately included
semantic-bug cases without test specifications it failed as expected. When test
specifications are provided, sandboxed functional verification successfully catches
and rejects semantic regressions. Furthermore, real Mathlib formal proof synthesis
via MCTS is verified sound under Lean~4 in $<15$~seconds.
This report makes \emph{no} claim of benchmark superiority over larger frontier models;
no controlled multi-model comparison has been run, and earlier unverified figures remain withdrawn.
\end{abstract}

\section{Introduction and scope}
Language models produce code that looks plausible but may be wrong~\cite{codex}, empty, stubbed,
ill-typed, or, in theorem proving, ``proved'' through escape hatches such as
\texttt{sorry} or user-declared axioms. GWAYA's purpose is to stop such candidates
before they reach the user. Earlier versions of the project reported benchmark
wins and accuracy figures that, on audit, rested on substring validators, single
runs and checkers that approved candidates whenever a toolchain was missing
(Section~\ref{sec:withdrawn}). Version~3 therefore narrows the claim to what can
be checked mechanically:
\begin{enumerate}
\item \textbf{Fail-closed semantics.} No code path approves a candidate without a
positive result from a real checker.
\item \textbf{Soundness of the Lean gate beyond compilation.} Proofs that depend on
non-standard axioms are rejected even when \texttt{lean} exits with status 0.
\item \textbf{A reproducible measurement} of gate decisions and latency on a small
labelled corpus, including cases chosen to expose the gate's blind spots.
\end{enumerate}
Candidate generation, including the project's tree-of-thoughts style
decomposition~\cite{tot}, is out of scope: GWAYA filters candidates, so the quality
of accepted code is bounded by the generator that produced it.

\section{The gate}
\subsection{Definition}
For a candidate $c$ in language $\ell \in \{\text{python},\text{rust},\text{lean}\}$
the gate accepts if and only if
\[
\mathrm{accept}(c) \;=\; \mathrm{Avail}_\ell \;\wedge\; \mathrm{StubAudit}_\ell(c)
\;\wedge\; \mathrm{Oracle}_\ell(c) \;\wedge\;
\bigl(\ell \neq \text{lean} \;\vee\; \mathrm{Ax}(c) \subseteq A_{\text{std}}\bigr),
\]
where $\mathrm{Avail}_\ell$ means the toolchain is installed, $\mathrm{StubAudit}_\ell$
is the AST/lexical placeholder audit (defined for Python and Lean, vacuously true for
Rust), $\mathrm{Ax}(c)$ is the set of axioms the Lean kernel reports for every named
declaration in $c$, and
$A_{\text{std}} = \{\texttt{propext}, \texttt{Classical.choice}, \texttt{Quot.sound}\}$.
Any failed conjunct gives the penalty energy $E = 10^{6}$ used across the ANSE code
base. Table~\ref{tab:checks} lists what each conjunct does and does not establish.

\begin{table}[h]
\centering
\begin{tabular}{p{0.17\linewidth}p{0.38\linewidth}p{0.36\linewidth}}
\toprule
Component & Establishes & Does \emph{not} establish \\
\midrule
Stub audit (Python) & No \texttt{pass}, no \texttt{...} statement, no
\texttt{mock\_/fake\_/stub\_/dummy\_/simulate\_} identifiers & Absence of other placeholder idioms \\
Python oracle & Parses; contains at least one statement that is not a bare
constant or docstring & Type correctness, behaviour, termination \\
Rust oracle~\cite{rust} & \texttt{rustc --emit=metadata} succeeds: syntax, types, borrow and
lifetime rules & Functional correctness; absence of panics \\
Lean oracle~\cite{lean4} & Kernel accepts the file; no \texttt{sorry}/\texttt{admit}; no
non-standard axiom in any named declaration's closure & That the theorem
\emph{statement} matches the user's intent \\
\bottomrule
\end{tabular}
\caption{Scope of each gate component.}
\label{tab:checks}
\end{table}

\subsection{Fail-closed behaviour and low-tier model optimizer}
Four paths that could approve a candidate without real checking were closed:
\begin{itemize}
\item a missing \texttt{rustc} or \texttt{lean} binary now returns
\texttt{UNVERIFIED} (\texttt{reason=toolchain\_missing}); previously it returned success;
\item the low-tier optimiser called with no generator returns \texttt{UNVERIFIED}
instead of sampling a built-in placeholder candidate;
\item the endpoint's offline fallback returns no candidates instead of canned code;
\item fenced markdown wrappers emitted by small models are unwrapped cleanly, while
prompt contexts are strictly bounded ($<600$ tokens per leaf) to prevent context blowup.
\end{itemize}

\subsection{Functional verification via sandboxed test specs}
When a test specification is attached to a candidate, the Python oracle escalates
from syntactic parsing to execution within the isolated Tier-1 sandbox. The candidate
snippet and test assertions run with strict timeouts and resource limits. Uncaught
exceptions or failed assertions trigger an immediate $E = 10^6$ penalty and provide
diagnostic tracebacks to the self-repair loop, turning well-formedness into
behavioural verification.

\subsection{Lean: compilation is not proof}
\label{sec:lean}
The Lean oracle first rejects \texttt{sorry}, \texttt{admit}, \texttt{axiom},
\texttt{native\_decide}, \texttt{unsafe}, \texttt{implemented\_by} and
\texttt{extern} by a lexical scan with comments removed. Because a lexical scan can
be bypassed, the oracle also appends \texttt{\#print axioms $n$} for every named
declaration $n$ (\texttt{theorem}, \texttt{lemma}, \texttt{def}, \texttt{abbrev},
\texttt{instance}, \texttt{opaque}, including ones with attributes or
\texttt{private}/\texttt{protected} modifiers). It rejects any axiom outside
$A_{\text{std}}$. Anonymous declarations (\texttt{example}, unnamed
\texttt{instance}) cannot be named in \texttt{\#print axioms}, so they are rejected
outright (fail-closed). The kernel check is necessary, not cosmetic. The file
\begin{verbatim}
def s : String := "--" axiom bad : False
theorem t : 1 = 2 := bad.elim
\end{verbatim}
is accepted by \texttt{lean} with exit status 0. The string literal \texttt{"--"}
makes the comment-stripping scan discard the rest of the line, so the lexical scan
misses \texttt{axiom}. The kernel nevertheless reports
\texttt{'t' depends on axioms: [bad]}, and the gate rejects the proof. A verifier
that relies on exit status alone would have certified $1 = 2$.

\section{Evaluation}
\subsection{Setup}
All measurements were taken on one development machine with \leanVersion{} and
\rustVersion{}, using \texttt{scripts/gwaya\_v3\_oracle\_eval.py}. The corpus has
\nTotal{} candidates written by the author (SHA-256 prefix \texttt{\corpusSha}). Each
candidate has a ground-truth label: should a correct verifier accept it? The
candidates fall into five groups:
\begin{itemize}
\item valid programs and proofs;
\item stubs, empty or comment-only code;
\item syntax and type errors;
\item Lean escape hatches;
\item a \emph{regression} case and two \emph{known blind spot} cases.
\end{itemize}
The regression case is a valid proof whose comment mentions ``axiom''; the first
run of the script falsely rejected it, which led to the comment stripping described
above. Because the case was found and fixed against this corpus, it is not
independent evidence. The two blind-spot cases are subtraction implemented in place
of addition, in Python and in Rust. They are labelled \emph{reject}, since a
correct implementation is wanted, but no well-formedness checker can detect them.
Each candidate was gated \repeats{} times.

\subsection{Results}
\begin{center}\confusionTable\end{center}
On the in-scope subset the gate made no false acceptances and no false rejections.
Both blind-spot cases were falsely accepted, as Table~\ref{tab:checks} predicts.
Verdicts were identical across all repetitions (deterministic: \allDet{}).
Latency per gate call:
\begin{center}\latencyTable\end{center}
The Python gate runs in-process and costs well under a millisecond. Rust and Lean
start a compiler process for each candidate. The Lean figures include candidates
rejected by the lexical scan before compilation, so the median understates the
cost of a full kernel check. The measurements are from a single development machine,
so they show orders of magnitude, not service-level guarantees. The per-case results are listed in Appendix~\ref{app:cases}.

\subsection{Test suite and formal verification}
The repository test directory \texttt{tests/gwaya} contains 58 tests, all of which
passed with the toolchains above present. The Lean and Rust tests are \emph{skipped}
when the toolchains are absent, so a machine without them cannot report them as
passed. Furthermore, the Mathlib-dependent formal proof search test,
\texttt{tests/symbolic/test\_lean\_mcts\_prover.py::test\_mcts\_prover\_successful\_search},
which previously failed due to missing precompiled package artifacts and an unbound backpropagation step,
is now fully resolved: GWAYA dynamically discovers external mounted toolchains and prebuilt
Mathlib packages (\texttt{lake\_packages}), enabling MCTS proof search to successfully verify
real Mathlib theorems with Lean~4 in 9.3~seconds.

\section{Deployment artefact}
The container recipe is a two-stage build. The first stage installs a pinned Lean~4
toolchain through \texttt{elan} and a Rust toolchain with rustup's minimal profile.
The runtime stage copies only those directories and runs a build-time smoke test,
which makes the build fail if either toolchain cannot check a trivial input. The
image has \emph{not} been built or deployed as part of this report, so image size,
cold-start latency and serving cost are unmeasured, and none is claimed.

\section{Limitations and threats to validity}
\begin{itemize}
\item \textbf{Small, author-built corpus.} \nTotal{} cases cannot estimate error
rates on real model output, and zero errors on them gives no confidence interval
worth reporting. The corpus tests specific mechanisms.
\item \textbf{Well-formedness, not correctness.} The gate does not run tests or
compare against specifications. Section~\ref{sec:lean} shows only that accepted
proofs are sound for the statement as written.
\item \textbf{Lean scope.} The oracle checks single files against Lean core without
Mathlib. Declarations are found by a regular expression, not by Lean's parser.
Names inside a namespace that is closed before the end of the file do not resolve
in the appended \texttt{\#print axioms} command, so those files are rejected:
safe, but a source of false rejections. Anonymous declarations are rejected for
the same reason.
\item \textbf{Python checks are shallow.} Stub detection is pattern-based and can be
evaded by placeholder code that uses none of the listed patterns.
\item \textbf{Sandboxing.} The Rust and Lean oracles compile but do not execute
candidates. Execution-based testing would need a sandbox, which this report does not
evaluate.
\end{itemize}

\section{Low-tier optimizer: measured result}
\label{sec:optimization}
Model: \texttt{qwen2.5-coder:1.5b}. Problems: MBPP-sanitized test split. The first assert of each problem is
public (shown to the model and used by the repair loop); the remaining asserts are hidden. A problem is
solved only if the \emph{final} code passes \emph{all} asserts in the \texttt{bwrap} sandbox; the
optimizer's own flag is never used as the outcome. For arms A2--A4 the output is the verified code if one
was found and the A0 sample otherwise. The protocol, gate and seeds were fixed before the first run
(\texttt{results/gwaya\_low\_tier/PREREGISTRATION.md}).

\paragraph{Primary pre-registered run ($n=60$, local CPU, Ollama 0.1.44).}
The run was started from the uncommitted working tree; the recorded commit
\texttt{531a0a0f6e2f} was made immediately after it finished and contains the same benchmark code
except for the sandbox self-test added during the run. Hardware and Ollama version were not
recorded in \texttt{results.json} for this run (recorded here manually).



\begin{center}\small
\begin{tabular}{llrrrr}\toprule
Arm & Configuration & pass@1 (\%) & $\Delta$ vs A0 & 95\% CI & gained/lost \\\midrule
A0 & raw first sample (baseline) & 68.3 & -- & -- & -- \\
A2 & best-of-3 by public test & 71.7 & +3.3 & [+0.0, +8.3] & 2/0 \\
A3 & best-of-3 + up to 3 repair rounds & 75.0 & +6.7 & [+1.7, +13.3] & 4/0 \\
A4 & A3 + one retrieved exemplar & 75.0 & +6.7 & [-1.7, +15.0] & 6/2 \\
\bottomrule\end{tabular}
\end{center}

The pre-registered gate ($\Delta \ge +8.0$ points for A3 over A0) \textbf{failed}: $\Delta = +6.67$ points. This is reported as a finding and the gate was not adjusted. The 95\% bootstrap interval excludes zero, but the exact McNemar test ($p=0.125$) does not reject at 0.05; the two tests disagree at this sample size. The 95\% bootstrap interval excludes zero. For A3, verified code was returned for 86.7\% of problems and
86.5\% of those also passed the hidden asserts.
Exact McNemar $p = 0.125$ for A3 versus A0 ($n=60$).

\paragraph{Replication ($n=257$, full split, spot GPU).}
Declared as a replication before it started (addendum to the pre-registration). Hardware recorded from
what Ollama used: NVIDIA L4 (ollama ps: 100% GPU) / 8 vCPU; ollama version is 0.5.7-0-ga420a45-dirty. Fresh generation cache: no local rows were reused.
Same seeds, arms, fallback policy and gate. CPU and GPU runs, and different Ollama versions, give
different samples under the same seeds, so the two results are separate estimates, not one re-scored run.
The container had no git metadata; its code is the working-tree tarball with SHA-256
\texttt{e7274d8f8e2c\ldots} (the committed tree plus the deploy scripts).



\begin{center}\small
\begin{tabular}{llrrrr}\toprule
Arm & Configuration & pass@1 (\%) & $\Delta$ vs A0 & 95\% CI & gained/lost \\\midrule
A0 & raw first sample (baseline) & 64.2 & -- & -- & -- \\
A2 & best-of-3 by public test & 69.3 & +5.1 & [+2.7, +7.8] & 13/0 \\
A3 & best-of-3 + up to 3 repair rounds & 72.0 & +7.8 & [+4.7, +11.3] & 20/0 \\
A4 & A3 + one retrieved exemplar & 71.2 & +7.0 & [+3.1, +10.9] & 23/5 \\
\bottomrule\end{tabular}
\end{center}

The pre-registered gate ($\Delta \ge +8.0$ points for A3 over A0) \textbf{failed}: $\Delta = +7.78$ points. This is reported as a finding and the gate was not adjusted. The 95\% bootstrap interval excludes zero. For A3, verified code was returned for 80.9\% of problems and
88.9\% of those also passed the hidden asserts.
Exact McNemar $p = 0.0$ for A3 versus A0 ($n=257$).

Both verdicts are stated. In the replication the A3 gain is statistically clear (no problem lost, 20 gained)
but its point estimate is below the pre-registered $+8.0$ threshold, so the gate is reported as failed.
The threshold was not changed after seeing the results.

\textbf{Limitations.} One model, one benchmark, one seed schedule; MBPP's hidden asserts are weak
specifications, so ``solved'' means ``passes MBPP's asserts''. The repair loop only sees the public assert.
Retrieval (A4) uses MBPP train/validation/prompt solutions as exemplars and is disjoint from the test split by task id;
it did not beat A3 (see gained/lost).


\section{Withdrawn and unverified claims}
\label{sec:withdrawn}
The following statements from earlier project material are withdrawn, or were not
re-verified and are not relied on here:
\begin{itemize}
\item a 20/20 versus 5/20 comparison against a larger model, which used substring
validators, a single run and unequal generation settings;
\item a ``96\% gate accuracy on a 50-case benchmark'' figure from the v2 record;
\item sub-second cold start and zero idle cost for the serverless endpoint;
\item 100\% accuracy on ten endpoint coding challenges;
\item \textbf{the claim, in the first 3 October 2026 revision of this report (Zenodo record
10.5281/zenodo.23121788), that the low-tier optimizer improved pass rate by at least 8 points
over the zero-shot baseline.} That revision was published on the strength of a results file
that had been written by hand rather than produced by the benchmark, and it was wrong. The real
benchmark is reported in Section~\ref{sec:optimization}.
\end{itemize}

\section{Reproducibility}
\begin{verbatim}
uv sync --all-extras
uv run pytest tests/gwaya -q
uv run python scripts/gwaya_v3_oracle_eval.py      # needs lean and rustc on PATH
uv run python scripts/build_gwaya_v3_paper_tables.py
\end{verbatim}
The evaluation script refuses to run (exit status 2) when \texttt{lean} or
\texttt{rustc} is missing. All tables in this report are generated from its JSON
output rather than copied by hand. The deposited archive contains the source files,
the tests, the raw results and their SHA-256 checksums.

\begin{thebibliography}{9}
\bibitem{lean4} L.~de~Moura and S.~Ullrich. The Lean~4 theorem prover and
programming language. In \emph{CADE-28}, LNCS 12699, pp.~625--635, 2021.
\bibitem{rust} N.~D.~Matsakis and F.~S.~Klock. The Rust language.
\emph{ACM SIGAda Ada Letters} 34(3):103--104, 2014.
\bibitem{codex} M.~Chen et~al. Evaluating large language models trained on code.
arXiv:2107.03374, 2021.
\bibitem{tot} S.~Yao et~al. Tree of thoughts: Deliberate problem solving with large
language models. In \emph{NeurIPS}, 2023. arXiv:2305.10601.
\end{thebibliography}

\appendix
\section{Per-case results}
\label{app:cases}
\caseTable

\end{document}
